\documentclass[11pt,a4paper]{article}
% Common preamble for the hint sheets (\input from each file)
\usepackage[utf8]{inputenc}
\usepackage[T1]{fontenc}
\usepackage{lmodern}
\usepackage{amsmath,amssymb,amsthm,mathtools}
\usepackage[a4paper,margin=2.3cm]{geometry}
\usepackage{enumitem}
\usepackage{xcolor}
\usepackage[hidelinks]{hyperref}
\renewcommand{\leq}{\leqslant}
\renewcommand{\le}{\leqslant}
\renewcommand{\geq}{\geqslant}
\renewcommand{\ge}{\geqslant}
\setlength{\parindent}{0pt}
\setlength{\parskip}{0.35em}
% \hintstitle{Sheet title}{Class work / Homework}
\newcommand{\hintstitle}[2]{\begin{center}{\LARGE\bfseries Hints: #1\par}\vspace{0.5em}{\large #2\par}\end{center}
\noindent\textit{The hints to each problem are cumulative. Read one, try again, and open the next only if you are still stuck.}\par\medskip\hrule\medskip}
\newcommand{\problem}[1]{\par\bigskip{\large\bfseries Problem #1}\par\nopagebreak\smallskip}
\newcommand{\subpart}[1]{\par\smallskip\textbf{(#1)}\par\nopagebreak}
\newcommand{\hint}[1]{\par\smallskip\noindent\textcolor{black!70}{\textbf{Hint #1.}}\enspace}
\newcommand{\note}[1]{\par\smallskip\noindent\textit{Note.} #1\par}
\providecommand{\lt}{<}
\providecommand{\gt}{>}
\providecommand{\tightlist}{\setlength{\itemsep}{0pt}\setlength{\parskip}{0pt}}
\emergencystretch=3em

\begin{document}
\hintstitle{Number theory -- 3}{Homework}

\problem{1}
\hint{1} For each $x$, the number of $y$ with $y^2=c$ is $1+\bigl(\frac cp\bigr)$. So $N=p+\sum_x\bigl(\frac{1-x^2}p\bigr)$.
\hint{2} Key lemma: $\sum_x\bigl(\frac{ax^2+bx+c}p\bigr)=-\bigl(\frac ap\bigr)$ if $b^2-4ac\ne0$, and $(p-1)\bigl(\frac ap\bigr)$ if $b^2-4ac=0$. Prove it by completing the square and reducing to $\sum_x\bigl(\frac{x^2-d}{p}\bigr)$, $d\ne0$. That sum equals $-1$: count the solutions of $x^2-y^2=d$, i.e.\ $(x-y)(x+y)=d$, in two ways.
\hint{3} Hence $N=p-\bigl(\frac{-1}p\bigr)$. For a general conic, complete squares to reach $au^2+bv^2=c$ and apply the lemma. A non-degenerate conic has $p-\bigl(\frac{-ab}p\bigr)$ affine points.

\problem{2}
\hint{1} By M\"obius inversion, $N_q=\sum_{d\mid q}\mu(d)\lfloor q/(4d)\rfloor$. For odd $e$, $\lfloor e/4\rfloor=\frac{e-2+\chi_4(e)}4$, where $\chi_4$ is the non-trivial character mod $4$.
\hint{2} For $q>1$ this gives $4N_q=\varphi(q)+\chi_4(q)\prod_{p\mid q}(1-\chi_4(p))$. The product is $0$ if some $p\equiv1\pmod 4$ divides $q$, and $2^{\omega(q)}$ otherwise.
\hint{3} $N_q$ is odd if and only if $4N_q\equiv4\pmod 8$. Since $2^{\omega(q)}\mid\varphi(q)$, two or more prime factors give an even $N_q$. For $q=p^k$, check $p\bmod 8$ and the parity of $k$: $p\equiv5$ (where $N_q=\varphi(q)/4$) and $p\equiv7$ give odd $N_q$, while $p\equiv1,3$ give even $N_q$.

\problem{3}
\hint{1} Thue: consider $u-av$ for $0\le u,v\le\lfloor\sqrt p\rfloor$. There are more than $p$ such pairs, so two coincide mod $p$; subtract. With $a^2\equiv-1$ this gives $p\mid x^2+y^2<2p$.
\hint{2} For $p\equiv1\pmod3$, $-3$ is a quadratic residue (reciprocity). Take $a^2\equiv-3$. Then $p\mid x^2+3y^2<4p$, so $x^2+3y^2\in\{p,2p,3p\}$.
\hint{3} $2p$ is impossible: $x,y$ would have the same parity, giving $4\mid x^2+3y^2$. If it is $3p$, then $3\mid x$ and $p=y^2+3(x/3)^2$.

\problem{4}
\hint{1} Squares are $0,1,4\pmod8$, so three of them never sum to $7\pmod 8$. If $4\mid x^2+y^2+z^2$, then all of $x,y,z$ are even. Descend.
\hint{2} By Euler's four-square identity it suffices to treat primes. For odd $p$, the sets $\{x^2\}$ and $\{-1-y^2\}$ each have $\frac{p+1}2$ elements mod $p$, so they meet: $x^2+y^2+1\equiv0$, i.e.\ $mp$ is a sum of four squares for some $m<p$.
\hint{3} Take $m$ minimal and suppose $m>1$. Reduce the $x_i$ to $y_i\equiv x_i\pmod m$ with $|y_i|\le m/2$; then $\sum y_i^2=mr$ with $r<m$. Multiply the two representations using Euler's identity: every term is divisible by $m$. Dividing by $m^2$ represents $rp$, a contradiction. (Treat the case where all $|y_i|=m/2$ separately.)

\problem{5}
\hint{1} Use multiplicativity in both arguments. The exponent $\frac{n-1}2$ is additive mod $2$ on odd numbers: $\frac{ab-1}2\equiv\frac{a-1}2+\frac{b-1}2\pmod 2$. So the reciprocity sign for $m,n$ is the product of the signs for all prime pairs $p\mid m$, $q\mid n$. The supplements follow in the same way.
\hint{2} Example: $\bigl(\frac2{15}\bigr)=\bigl(\frac23\bigr)\bigl(\frac25\bigr)=1$, yet $2$ is not a square mod $3$, hence not mod $15$.
\hint{3} Solovay--Strassen: test whether $a^{(n-1)/2}\equiv\bigl(\frac an\bigr)\pmod n$. The Jacobi symbol is computed fast by reciprocity. For odd composite $n$ the $a$ passing the test form a proper subgroup of $(\mathbb Z/n)^*$, so a random $a$ fails with probability $\ge\frac12$.

\problem{6}
\note{The equation is $\sqrt[3]{7x^2-13xy+7y^2}=|x-y|+1$.}
\hint{1} $x=y$ gives $(1,1)$. Otherwise let $x>y$ and $n=x-y$, and cube both sides. You get $y^2+ny=n^3-4n^2+3n+1$.
\hint{2} Complete the square: $(2y+n)^2=(n-2)^2(4n+1)$. So $4n+1$ is an odd square, $n=m^2+m$ with $m\ge2$.
\hint{3} Then $2y+m^2+m=(m-1)(2m^2+5m+2)$. Answer: $(1,1)$ and $\{x,y\}=\{m^3+2m^2-m-1,\ m^3+m^2-2m-1\}$, $m\ge2$.

\problem{7}
\hint{1} $p=2$: all $x+y=2^k$ work. For odd $p$, let $x^{p-1}+y=p^a$ and $x+y^{p-1}=p^b$ with $x\le y$, so $a\le b$. Substitute $y=p^a-x^{p-1}$ and work mod $p^a$.
\hint{2} This gives $p^a\mid x^{(p-1)^2}+x$. Show $p\nmid x$, so $p^a\mid x^{p(p-2)}+1$, and $p\mid x+1$.
\hint{3} By LTE, $v_p(x^{p(p-2)}+1)=v_p(x+1)+1=:r+1$, so $a\le r+1$, while $p^r\le x+1\le p^a$. This forces $x=p-1$ and $a=2$. Then $(p-1)^{p-1}<p^2$ forces $p=3$. Answer: $(3,2,5)$, $(3,5,2)$, and $p=2$ with $x+y$ a power of $2$.

\problem{8}
\hint{1} Five suffice: $x^2+7=x^2+2^2+1^2+1^2+1^2$. With four squares, each $f_i=a_ix+b_i$ is linear and $\sum a_i^2=1$, $\sum a_ib_i=0$, $\sum b_i^2=7$.
\hint{2} So $b$ lies in the orthogonal complement of a rational unit vector $a$ in $\mathbb Q^4$. By Witt cancellation, the form $\sum x_i^2$ restricted to $a^\perp$ is $\mathbb Q$-equivalent to $x^2+y^2+z^2$. Hence $7$ would be a sum of three rational squares.
\hint{3} Clearing denominators, $7m^2$ would be a sum of three integer squares. But $7m^2$ has the form $4^k(8l+7)$ (Problem~4). (The official solution does the descent mod $8$ by hand.) Answer: $n=5$.

\problem{9}
\hint{1} For $p\equiv1\pmod4$, $\bigl(\frac{p-a}p\bigr)=\bigl(\frac ap\bigr)$. Pair $a$ with $p-a$: each pair contributes $p\bigl(\frac ap\bigr)$.
\hint{2} The sum equals $p\sum_{a<p/2}\bigl(\frac ap\bigr)$. That inner sum is half of $\sum_{a=1}^{p-1}\bigl(\frac ap\bigr)=0$, so the whole sum is $0$.
\hint{3} For $p\equiv3\pmod4$, $p>3$, Dirichlet's class number formula gives $h(-p)=-\frac1p\sum_{a=1}^{p-1}\bigl(\frac ap\bigr)a$. For example, $p=7$ gives $-7$, i.e.\ $h=1$; $p=23$ gives $-69$, i.e.\ $h=3$.

\problem{10}
\hint{1} $g=\sum_a\bigl(\frac ap\bigr)\zeta^a$ with $\zeta=e^{2\pi i/p}$. Compute $g\bar g$ or $g^2=\sum_{a,b}\bigl(\frac{ab}p\bigr)\zeta^{a+b}$: substitute $b=at$ and sum over $a$ first. You get $g^2=\bigl(\frac{-1}p\bigr)p=:p^*$.
\hint{2} In $\mathbb Z[\zeta]$, $g^q\equiv\sum\bigl(\frac ap\bigr)\zeta^{qa}=\bigl(\frac qp\bigr)g\pmod q$ (Frobenius, using $q\mid\binom qj$).
\hint{3} Also $g^q=g\,(p^*)^{(q-1)/2}\equiv\bigl(\frac{p^*}q\bigr)g$. Multiply by $g$ and cancel the unit $p^*$ mod $q$: $\bigl(\frac qp\bigr)=\bigl(\frac{p^*}q\bigr)$, which is reciprocity.
\end{document}
